\begin{exsol}{2}{\MR{} Food for Thought 1.2}{18}{Foundations}{A1}
\asked For an army in battle, identify which decisions are taken centrally and which are left to
field commanders, state the principles behind the division, and judge whether the
coordination/motivation split helps in analysing military organization.

\method Coordination versus motivation (\MR{} Ch.~1, p.~16; Ch.~2, pp.~25--27). Authority goes to
those with the information, and delegation is valuable only with aligned objectives (p.~17). Design
attributes, synchronization and assignment problems (Ch.~4, pp.~91--92). Urgency plus central
knowledge of the optimum favours central control (p.~92). Routines (p.~92). Central decisions set
the parameters of local ones (pp.~114--116).

\solution
\subparagraph{What is decided centrally.} The objective of the operation and the main line of
effort; the \emph{timing} of the attack; the boundaries between units; the commitment of reserves;
the allocation of artillery, air support and supplies among units; and the rules of engagement.
Each is a design variable (p.~91). Timing is a \emph{synchronization} problem: an assault on two
axes, or infantry following a barrage, works only if everyone moves together. A unit that attacks
early or late is destroyed piecemeal, which is the coxswain's problem of p.~91 with far higher
stakes. Unit boundaries and fire support are \emph{assignment} problems (pp.~91--92): each task
needs exactly one unit, and the errors are a gap in the line (no one assigned) or fire on one's
own troops (two assigned). As with the ambulance dispatcher, a central decision prevents
duplication and delay, and a small error in the overall plan costs far less than a failure of fit.
The centre also collects intelligence from many sources (reconnaissance, other sectors), which no
field commander has. This is an innovation attribute (p.~92): the needed information is not
available to any one unit.

\subparagraph{What is left to field commanders.} How to take the objective given the terrain;
fire and movement; when to exploit an opening or fall back; how to react when the enemy does
something unexpected. The information is local, perishable and observed only on the spot.
Communication is slow and is destroyed by the battle itself, so routing it upward would make the
decision both late and uninformed --- the HBC's failure of over-centralization (p.~26) in its
starkest form. Hence a familiar military principle: the centre issues the \emph{intent} (what and
why, by when) and leaves the \emph{method} (how) to the subordinate. This is David Packard's rule
quoted on p.~116: agree on the objective, then ``turn everybody loose''. Drills and standard
procedures are routines (p.~92) that let units act coherently even when contact with the centre
is lost.

\subparagraph{Principles.} (1) Centralize what has design attributes: timing, boundaries,
assignments, the overall plan. (2) Decentralize what depends on local, fast-changing information.
(3) The more urgent and disrupted the communication, the more must be decided below, within
parameters set above (p.~114).

\subparagraph{Coordination and motivation.} The split is useful, and motivation in an army is
extreme: the organization asks people to risk their lives, so individual interest (survival)
and the organization's goal (winning) diverge sharply. Output is produced by teams and is hard to
measure, so pay for performance is nearly useless. Armies rely instead on selection and training,
discipline and punishment, unit cohesion and loyalty, honours, and promotion based on careers. But
the two problems interact. Delegation to field commanders pays only if they share the objective
(p.~17): an army that cannot trust its officers' competence or loyalty must centralize more and
accept slower, less informed decisions. Motivation also shapes coordination through reporting.
Commanders who are judged on success have reasons to overstate progress, which corrupts the
information the centre plans with.

\interp The army shows the course's two principles at their limit. Central control of design
variables buys fit, local control buys speed and information, and the boundary between them is set
by the cost of misfit, the speed of change and the trust that incentives and selection create.

\verify Cross-check against the opposite extreme. A fully centralized army, with every tactical
move approved above, fails on information and speed. A fully decentralized one, with each unit
choosing its own objective and timing, fails on synchronization and assignment. Both fail in the
predictable ways of p.~93 (central commands: quick but ``insensitive to local knowledge'';
decentralized: duplication and delay).

\answer{The centre sets objectives, timing, unit boundaries, reserves, fire support and supply
allocation --- design variables with synchronization and assignment attributes, where misfits are
catastrophic. Field commanders choose the methods, using local, perishable information within the
centre's intent. The coordination/motivation split is useful, but the two interact: delegation is
viable only with trained, loyal and motivated commanders, and incentive problems distort the
information that coordination needs.}
\end{exsol}
